% Solutions/hints for ch24 exercises (assembled into Appendix C)

\begin{solution}{exr:ch24-horizon}
(a) The stop form of \cref{eq:ch24-tauh} gives $\tau_h \ge \dt + v_{\max}/a_{\max} = 0.05 + 1 = 1.05$\,s, the reversal form $\tau_h \ge 0.05 + 2 = 2.05$\,s, and the hand-over condition $\tau_h \ge \ttc = 2$\,s. If a stop is enough, $\tau_h = 2$\,s satisfies everything; if the local layer may have to reverse the velocity, $\tau_h = 2.05$\,s is the smallest choice, and rounding to a multiple of $\dt$ gives $2.05$\,s exactly. (b) The intruder must be tracked before it enters the horizon: the range must cover the closing distance during the warm-up and the horizon plus the inflated separation, $(v_{\mathrm{cruise}} + v_B)(\tau_h + n_{\mathrm{warm}}\dt) + R_{\mathrm{safe}} + \kappa_p\sigma(\tau_h) = (1.5 + 3)(2 + 20 \cdot 0.05) + 1 + 0.5 = 15$\,m, the closing speed being the cruise speed because the drone is Nominal while it waits for the trigger; a drone that may already be flying at $v_{\max} = 2$\,m/s when the intruder appears needs $(2 + 3) \cdot 3 + 1.5 = 16.5$\,m. (c) With $\tau_h < \ttc$ the executive hands control to \orca only when the current relative velocity already lies inside the truncated velocity obstacle $\VO^{\ttc}$, so the first half-plane demands a large change of velocity at once, which a drone with an acceleration limit may not be able to execute and which makes the program more likely to be infeasible; with $\tau_h \ge \ttc$ the first correction is small. With $\tau_h$ far beyond the look-ahead at which $\kappa_p\sigma(s)$ reaches the lane spacing, the inflated disk of an intruder in the \emph{next} lane covers the drone's own lane, the trigger fires for intruders that would never come close, drones leave their plans needlessly, and every needless deviation costs reconnections and re-checks; the trigger becomes a false-alarm generator rather than a safety device.
\end{solution}

\begin{solution}{exr:ch24-reconnect-bound}
(a) Line~\ref{alg:ch24-loop:shift} of \cref{alg:ch24-loop} replaces $\pi_i$ by its tail from index $j$ and sets $t_0^i \gets t_0^i + j + \Delta$, so the waypoint $\pi_i[j]$ is now scheduled at grid time $t_0^i + j + \Delta$ and every later waypoint, including the goal, at exactly $\Delta$ steps later than before; the arrival time of $a_i$ rises by $\Delta$. No other path changes, so $\sumcost$ rises by exactly $\Delta$ and $\makespan$ by at most $\Delta$ (by less if another drone arrives later than $a_i$), provided the re-check finds no conflict, because a repair would change further paths. (b) With $\Delta = 0$ line~\ref{alg:ch24-loop:shift} is skipped and $(\pi_i, t_0^i)$ is unchanged, hence the reservation that the other drones see is unchanged; the plan in force was conflict-free before the deviation (the invariant maintained by \cref{thm:ch24-nominal,thm:ch24-replan}), so it is conflict-free after. The only new motion is the straight segment, which line~\ref{alg:ch24-reconnect:mate} of \cref{alg:ch24-reconnect} has verified against the teammates' intended positions and line~\ref{alg:ch24-reconnect:pred} against the prediction; physical separation on the segment is therefore assured without any re-check of the plans. (c) Build the simulation, step it until $t = 4.9$\,s, take C's nominal path from index $5$ with $t_0 = 6$ and B's nominal path with $t_0 = 0$, and call \code{first\_conflict} on the pair from $t = 4$: C occupies $(6,2)$ at $t = 7$ and so does B, giving $\langle B, C, (6,2), 7\rangle$. The teammate test of the segment checks only the interval of the segment itself, $t \in [4.9, 6]$, during which B is still at $(6,3)$; the conflict arises at $t = 7$, after C has rejoined its plan, and it is a property of the \emph{shifted plan}, which no line of \cref{alg:ch24-reconnect} examines. That is precisely why every delayed reconnection is followed by the re-check.
\end{solution}

\begin{solution}{exr:ch24-offsets}
Padding extends a path backward to its first cell, so a detector started at grid time $0$ places A at $(6,5)$ for all $t \le 9$; B's nominal path visits $(6,5)$ at $t = 3$, and the detector reports $\langle A, B, (6,5), 3\rangle$, a conflict between a drone that did not exist at $(6,5)$ at that time and one that flew through it six seconds before the replan. Started at $\lfloor t \rfloor = 9$, the detector compares A at $(6,5)$ for $t = 9, 10$, $(6,6)$ for $t = 11$, and so on along row~6, with B at $(6,1)$ at $t = 9$ and parked at $(6,0)$ from $t = 10$ (its repaired path) and with C along row~1 and row~2: no cell is shared at any time, and no swap occurs, so nothing is reported. Start times are read on line~\ref{alg:ch24-reconnect:j0} of \cref{alg:ch24-reconnect} ($j_0$ from $t/\Delta T - t_0^i$) and in the arrival time $t_j = (t_0^i + j + \Delta)\Delta T$; written on line~\ref{alg:ch24-loop:shift} of \cref{alg:ch24-loop} ($t_0^i \gets t_0^i + j + \Delta$); read on line~\ref{alg:ch24-replan:res} of \cref{alg:ch24-replan} (the reservation of every $(\pi_j, t_0^j)$) and on line~\ref{alg:ch24-replan:detect} (the detector); and written on line~\ref{alg:ch24-replan:install} ($t_0^i \gets t_{\mathrm{start}} - 1$) and after the local \cbs ($t_0^i \gets t_0^j \gets t_{\mathrm{start}}$). Every one of them is an absolute grid time, and $t_{\mathrm{start}}$ itself is computed from the absolute clock $t$.
\end{solution}

\begin{solution}{exr:ch24-particles}
The particle trigger fires at look-ahead $s$ when $\sum \set{w^{(n)} : \norm{\tilde{\pos}_i(s) - \pos^{(n)}_h} < R_{\mathrm{safe}}} > 1 - p$; no covariance is formed. On a bimodal prediction (half the weight on each side of a building) the Gaussian summary puts its mean in the gap between the modes and gets a covariance whose largest eigenvalue is roughly half the distance between them, so the inflated disk $R_{\mathrm{safe}} + \kappa_p\sigma(h)$ covers the empty lane \emph{between} the modes: the Gaussian trigger fires early, on a place the intruder will almost surely not be, and the replanner blocks that lane as well. The particle trigger fires only when one mode carries more than $1 - p$ of the weight within $R_{\mathrm{safe}}$, that is when the drone's own lane is the threatened one; it costs $\bigO{N}$ per look-ahead instead of $\bigO{1}$, and its threshold is a weight, not a count, so unequal weights are handled without resampling.
\end{solution}

\begin{solution}{exr:ch24-states}
(a) With a finite cost $c_{\mathcal{B}}$ the search may cross the tube, so the statement ``the replanned path is clear of the $p$-disc at every step'' is lost; what survives is \cref{thm:ch24-replan}, which is about conflicts between \emph{plans}, and completeness of the replanner, which blocking can destroy in a corridor. Make $c_{\mathcal{B}}$ depend on how deep the cell lies inside the inflated disk, for instance $c_{\mathcal{B}} = 1 + c_0\,(d^{\mathrm{safe}}_h - \norm{\vect{c}(v) - \hat{\pos}_{T+h}})/(\kappa_p\sigma(h))$, so that a sharp prediction produces a narrow, expensive tube and a vague one a wide, cheap one; a constant cost ignores the covariance and makes the same detour in both cases. Expect A to keep the row~6 detour while $c_{\mathcal{B}}$ exceeds the seven extra steps it costs and to cut through the tube once it is cheaper. (b) Emergency is entered from any state when the local program has been infeasible for $m$ consecutive cycles, brakes to hover (or climbs, if the third dimension is available), and is tested before every other transition of \cref{alg:ch24-loop}; it leaves to Avoiding when the program is feasible again for $n_{\mathrm{clear}}$ cycles, to Replanning when the drone is at rest and a path exists, and never directly to Nominal, because the plan in force has not been re-checked.
\end{solution}

\begin{solution}{exr:ch24-formation}
(a) C is a follower with $\vect{d}_{AC} = (0,-2)$, so while C was Nominal \cref{eq:ch24-formation} added $k_F(\pos_A + \vect{d}_{AC} - \pos_C)$ to its preferred velocity: when A climbed away from the intruder, C's slot climbed with it, C's \emph{intended} positions moved toward the inflated tube, and the trigger of \cref{def:ch24-inside}, which tests intended positions and not the current geometry, fired at $t = 4.3$\,s although the intruder never came closer to C than $2.25$\,m. Had the term stayed active while C was Avoiding, it would have pulled C back toward a slot that lies in the maneuver A is making, the \orca half-plane would have pushed it out again, and the two would have fought at $10$\,Hz; dropping the term while Avoiding is what \cref{sec:ch24-constraints} prescribes. (b) Expect the peak of $e_F$ to fall and its recovery to shorten as $k_F$ grows, at the price of a larger pull into the tube and a smaller minimum separation; ramping the gain over a second removes the step in $\vel^{\mathrm{pref}}$ at the entry to Reconnecting without changing the recovery time much. (c) When the \emph{leader} avoids, freeze the formation reference at the leader's nominal reference position $\pos^{\mathrm{ref}}_\ell(t)$ instead of its actual position (or hand leadership to a follower that is not in conflict) so that the whole formation does not follow one drone into its maneuver; restore the actual position when the leader is Nominal again.
\end{solution}

\begin{solution}{exr:ch24-dwa}
The rollouts need, per rollout step $h$, the predicted mean $\hat{\pos}_{T+h}$ and the inflated radius $R_{\mathrm{safe}} + \kappa_p\sigma(h)$ (and the teammates' intended positions with $R_{\mathrm{mate}}$), so that the clearance term of \cref{ch:ch14} is evaluated against a moving disk rather than a static obstacle. The heading term points at $\pos^{\mathrm{ref}}_i(t + t_{\mathrm{la}})$ while Nominal, at the reconnection target while Reconnecting, and at the same pursuit target with the formation term dropped while Avoiding, which keeps DWA's objective aligned with \cref{eq:ch24-vpref}. \orca's feasibility flag is replaced by ``the admissible window is empty'': no sampled $(v,\omega)$ has positive clearance over the rollout, which is the same escalation to Replanning on line~\ref{alg:ch24-loop:bound}. With a double integrator the dynamic window enforces $a_{\max}$ by construction, so the reversal form of \cref{eq:ch24-tauh} becomes the binding constraint on $\tau_h$; expect a slightly larger minimum separation (DWA brakes rather than sidesteps), a longer time in Avoiding and, because the drift bound is reached later, no more replans than with \orca.
\end{solution}
